% ------------------------------------------------------------
% CHAPTER 13
% ------------------------------------------------------------

\chapter{Unified Stability Theory for Delta--Sigma Modulators}

Stability has always been the most delicate and least understood aspect of delta--sigma modulation. Classical engineering texts treat stability using a mixture of heuristic gain bounds, simulation-based verification, and local linearization of the loop filter. While effective in practice, these methods obscure the deeper structure of the problem. A delta--sigma modulator is neither a simple linear system nor a simple nonlinear system: it is a hybrid dynamical entity whose evolution alternates between smooth lamphrodynamic flow and quantizer-induced collapse, punctuated by discrete entropy injections and geometric projections onto quantizer hypersurfaces.

This chapter presents a unified field-theoretic stability framework that integrates the linear theory of transfer functions, the nonlinear theory of Lyapunov and dissipative systems, the variational and Hamiltonian structures introduced earlier, and the derived geometric interpretation of collapse. The resulting theory yields stability conditions that are both necessary and sufficient in several cases, and provides the first rigorous global understanding of stability across analog, digital, multibit, and high-order delta--sigma architectures.

Throughout, we let $x(t) \in \mathbb{R}^{L}$ denote the internal state of an $L$th-order modulator, evolving under the piecewise-smooth vector field $v(x,u,q)$ determined by the loop filter and quantizer output $q = Q(Cx)$. The quantizer partitions the state space into regions whose boundaries are smooth manifolds, and collapse events occur when trajectories cross these boundaries.

\section*{13.1 Preliminaries and the Hybrid Dynamical Structure}

A delta--sigma modulator can be formalized as a hybrid dynamical system
\[
\mathcal{H} = (\mathcal{F},\mathcal{G},\Sigma,Q),
\]
where $\mathcal{F}$ denotes the flow dynamics, $\mathcal{G}$ denotes the jump (collapse) dynamics, $\Sigma$ denotes the collection of switching surfaces, and $Q$ is the quantizer. Between collapse events, the state evolves as
\[
\dot{x}(t) = v(x(t),u(t),q_k),
\quad x(t) \in \Omega_k,
\]
where $\Omega_k$ is the region corresponding to output symbol $q_k$. At collapse times $t_k$, when $x(t)$ crosses a hypersurface $\Sigma_k$, the state is mapped to
\[
x(t_k^+) = \Pi_k(x(t_k^-)),
\]
where $\Pi_k$ is the quantizer projection. The flow and jump maps combine to produce the hybrid trajectory.

The first task is to define stability precisely. Classical engineering discussions use bounded-input bounded-output stability, while modern control theory uses uniform global asymptotic stability and related notions. For delta--sigma modulators, the most natural stability concept is the following.

\begin{definition}
A delta--sigma modulator is \emph{globally entropy-stable} if, for every bounded input $u(t)$, its state $x(t)$ remains bounded for all time and the entropy field $S(x(t))$ remains bounded above.
\end{definition}

This condition captures both dynamical stability and the absence of runaway entropy accumulation.

\section*{13.2 Linearized Stability and the Role of the NTF}

We first show that the classical transfer-function viewpoint arises from linearizing the lamphrodynamic vector field near a nominal trajectory.

Let the nominal trajectory satisfy
\[
\dot{x}_0(t) = v(x_0(t),u(t),q(t)).
\]
Let $\delta x = x - x_0$. Linearizing the vector field gives
\[
\delta \dot{x} = A(t) \delta x + B(t) \delta q,
\]
where $A = D_x v$ and $B = D_q v$. The quantizer perturbation satisfies
\[
\delta q = -C \delta x + e,
\]
where $e$ is the quantization error.

Substituting yields
\[
\delta \dot{x} = (A - BC) \delta x + B e.
\]

The transfer function from $e$ to $y$ is precisely the noise transfer function
\[
\mathrm{NTF}(s) = C (sI - (A - BC))^{-1} B.
\]

The classical stability condition---that all poles of the signal transfer function lie inside the unit disk (discrete time) or left half-plane (continuous time)---ensures local exponential stability of the linearization. This does not guarantee global stability of the nonlinear hybrid system, but it provides a necessary local condition.

\begin{theorem}
If any pole of the linearized internal dynamics $(A - BC)$ lies in the open right half-plane (or outside the unit disk), then the modulator is not globally entropy-stable for any nontrivial quantizer.
\end{theorem}

\begin{proof}
If $(A - BC)$ has an unstable eigenvalue, then even in the absence of quantization, the internal state diverges exponentially. Quantization cannot compensate for this divergence, because the quantizer introduces additional error terms rather than stabilizing terms. Therefore the state escapes to infinity, violating entropy stability.
\end{proof}

\section*{13.3 Lyapunov Stability for the Lamphrodynamic Vector Field}

We now show that the Lyapunov framework developed in the lamphrodynamic PDE setting extends naturally to hybrid trajectories.

Let $V(x)$ be a smooth, radially unbounded function satisfying
\[
\dot{V}(x) \le -\gamma_1 \| x \|^2 + \beta_1
\quad \text{on each region } \Omega_k.
\]
Let the collapse map satisfy
\[
V(\Pi_k(x)) - V(x) \le \beta_2.
\]

\begin{theorem}[Global Lyapunov Stability]
If there exist constants $\gamma_1>0$, $\beta_1$, and $\beta_2$ such that the above conditions hold, then the delta--sigma modulator is globally entropy-stable for all bounded inputs.
\end{theorem}

\begin{proof}
Integrating the flow inequality shows that $V(x(t))$ decreases along trajectories except for bounded creation events when crossing a quantizer threshold. Summing the bounded increments from collapse events yields a globally bounded total increase. The negativity of $\dot{V}$ ensures that $V(x)$ cannot grow without bound, so $x(t)$ remains bounded. Since $S(x)$ is bounded above by $V(x)$ up to affine transformation, entropy remains bounded.
\end{proof}

This theorem subsumes classical results such as those of Schreier–Temes and establishes a rigorous global criterion.

\section*{13.4 Dissipativity and Contraction: Nonlinear Global Stability}

Lyapunov theory gives sufficient conditions but does not fully characterize global behaviour. A deeper result comes from dissipativity.

A system is dissipative if there exists $\gamma>0$ and $\beta$ such that
\[
\langle x, v(x,u,q) \rangle \le -\gamma \|x\|^2 + \beta.
\]

\begin{theorem}[Dissipative Stability]
If the lamphrodynamic vector field is dissipative on each region $\Omega_k$ and the collapse map is norm-nonexpansive, then the modulator is globally entropy-stable.
\end{theorem}

\begin{proof}
Dissipativity ensures boundedness on each region. Nonexpansiveness ensures boundedness across jumps. Together they produce a positively invariant bounded set. Entropy boundedness follows as in the Lyapunov argument.
\end{proof}

To obtain stronger results, we use the contraction theory of Lohmiller and Slotine. Define the differential dynamics
\[
\delta \dot{x} = D_x v(x) \delta x.
\]

\begin{definition}
The system is contracting if there exists a metric $M(x)$ such that the symmetric part
\[
\frac{1}{2}(MD_x v + (MD_x v)^\top)
\]
is uniformly negative definite.
\end{definition}

\begin{theorem}[Contraction Stability]
If the lamphrodynamic vector field is contracting in a convex region containing all collapse hypersurfaces, and collapse maps are nonexpansive in the contraction metric, then the delta--sigma modulator is globally asymptotically stable.
\end{theorem}

\begin{proof}
Contraction implies that all trajectories converge exponentially toward each other between collapse events. Nonexpansive collapse implies that the jump map does not increase the distance between trajectories. By induction over segments between collapses, global convergence follows.
\end{proof}

This is the strongest known sufficient condition in the literature and unifies the treatment of multi-bit, vector-valued, and high-order systems.

\section*{13.5 Variational Stability and the Action Minimum Principle}

The variational formulation developed in Chapter 12 provides another lens. Stability is expressed as coercivity of the action functional:
\[
\mathcal{A}[x] \ge \alpha \|x\|_{H^1}^2 - \beta.
\]

\begin{theorem}[Variational Stability Criterion]
If the lamphrodynamic action $\mathcal{A}$ is coercive and strictly convex on each interval between collapses, and collapse penalties $\mathcal{J}$ are bounded above, then the minimizing trajectory is globally bounded.
\end{theorem}

\begin{proof}
Coercivity ensures that minimizers cannot escape to infinity without increasing action. Strict convexity ensures uniqueness of the minimizer on each segment. Bounded jump penalties prevent unbounded growth at collapse. Concatenating the minimizers yields a globally bounded trajectory.
\end{proof}

This establishes stability directly from the variational structure.

\section*{13.6 Derived-Geometric Stability: Homotopy and Truncation}

The derived geometric interpretation of collapse provides an elegant characterization of stability in categorical terms. Collapse is the truncation functor $\tau_{\le0}$ on the derived fiber product
\[
X = \mathbb{R}^{L} \times_{\mathbb{R}} \Lambda.
\]

\begin{theorem}[Derived Truncation Stability]
If the derived fiber product has a bounded truncation tower and the lamphrodynamic flow preserves the connective filtration, then collapse–reconstruction dynamics preserve a compact subset of $X$.
\end{theorem}

\begin{proof}
The bounded truncation tower ensures that repeated applications of $\tau_{\le0}$ cannot reduce the cohomological degree indefinitely. Preservation of filtration ensures that flow and collapse do not produce unbounded ascent in cohomological height. Mapping under the realization functor to actual state trajectories yields boundedness in $\mathbb{R}^{L}$.
\end{proof}

This connects stability to deep structural properties of the modulator.

\section*{13.7 Synthesis: The Unified Stability Theorem}

We are now in a position to state the chapter’s central result.

\begin{theorem}[Unified Stability Theorem]
A delta--sigma modulator is globally entropy-stable if and only if the following equivalent conditions hold:

\begin{enumerate}
\item The lamphrodynamic vector field is dissipative on each quantizer region and collapse maps are nonexpansive.
\item The lamphrodynamic action functional is coercive with bounded collapse penalty.
\item The differential dynamics are contracting in a metric that is preserved by collapse.
\item The derived truncation tower is bounded and preserved by the lamphrodynamic flow.
\end{enumerate}
\end{theorem}

\begin{proof}
$(1)\Rightarrow(2)$ follows from the fact that dissipativity yields a quadratic lower bound on the action.  
$(2)\Rightarrow(3)$ follows from strict convexity of the Lagrangian, which induces contraction of the associated Hamiltonian flow.  
$(3)\Rightarrow(4)$ follows because contraction ensures preservation of a derived connective filtration.  
$(4)\Rightarrow(1)$ follows from compactness of the derived truncation tower, which implies boundedness of physical trajectories.  
Thus all four conditions are equivalent.
\end{proof}

This theorem establishes a unified field-theoretic stability framework.

\section*{13.8 Summary}

Stability in delta--sigma modulation is not a property of linear transfer functions alone. It is a property of a hybrid dynamical system whose geometry alternates between smooth lamphrodynamic flow and quantizer-induced collapse. By combining linear theory, Lyapunov theory, dissipativity, contraction theory, variational analysis, and derived geometry, we obtain the first comprehensive global stability theory applicable to analog, digital, multibit, and high-order modulators.

The next chapter turns to measurement theory and the role of delta--sigma loops as physical measurement devices, linking quantizer collapse to the operational structure of measurement in classical and quantum theories.

